\chapter{Forward--Backward Stochastic Differential Equations}

We study fully coupled FBSDEs of the form
\[
\left\{
\begin{aligned}
  X_t
  &=x+\int_0^t b_s(X_s,Y_s,Z_s)\,\mathrm ds
      +\int_0^t\sigma_s(X_s,Y_s,Z_s)\,\mathrm dB_s,\\
  Y_t
  &=g(X_T)+\int_t^T f_s(X_s,Y_s,Z_s)\,\mathrm ds
      -\int_t^T Z_s\,\mathrm dB_s.
\end{aligned}
\right.
\]
The processes $X,Y,Z$ take values in
$\mathbb R^{d_1}$, $\mathbb R^{d_2}$, and
$\mathbb R^{d_2\times d}$, respectively.  We write
$\Theta=(X,Y,Z)$ and
\[
  \|\Theta\|^2
  :=\mathbb E\!\left[
    \sup_{0\le t\le T}|X_t|^2
    +\sup_{0\le t\le T}|Y_t|^2
    +\int_0^T|Z_t|^2\,\mathrm dt
  \right].
\]

Unless stated otherwise, assume:
\begin{itemize}
  \item[(F1)] $\mathbb F=\mathbb F^B$ is the usual augmented Brownian
  filtration;
  \item[(F2)] with $\varphi_t^0:=\varphi_t(0,0,0)$ for
  $\varphi=b,\sigma,f$,
  \[
  \begin{gathered}
    g(0)\in L^2(\mathcal F_T;\mathbb R^{d_2}),\\
    \mathbb E\!\left[
      \left(\int_0^T(|b_t^0|+|f_t^0|)\,\mathrm dt\right)^2
      +\int_0^T|\sigma_t^0|^2\,\mathrm dt
    \right]<\infty;
  \end{gathered}
  \]
  \item[(F3)] $b,\sigma,f$, and $g$ are uniformly Lipschitz in their
  spatial variables.
\end{itemize}

\section{Local well-posedness by a contraction argument}

Let $L_{\sigma,z}$ denote the Lipschitz constant of $\sigma$ in $z$ and
let $L_g$ be the Lipschitz constant of $g$.

\begin{thm}[Small-time well-posedness]
Assume \textup{(F1)}--\textup{(F3)} and
\[
  L_{\sigma,z}L_g<1.
\]
There are constants $\delta_0>0$ and $C<\infty$, depending only on the
Lipschitz constants, the dimensions, and the product
$L_{\sigma,z}L_g$, such that, whenever $T\le\delta_0$, the FBSDE has a
unique solution and
\[
  \|\Theta\|^2\le C(|x|^2+I_0^2),
\]
where
\[
  I_0^2:=\mathbb E\!\left[
    \left(\int_0^T(|b_t^0|+|f_t^0|)\,\mathrm dt\right)^2
    +\int_0^T|\sigma_t^0|^2\,\mathrm dt
    +|g(0)|^2
  \right].
\]
\end{thm}

\begin{proof}
We present the scalar proof; the multidimensional proof differs only in
notation.  For progressively measurable input processes $(\mathbf y,
\mathbf z)$, first solve
\[
\begin{aligned}
  X_t^{\mathbf y,\mathbf z}
  &=x+\int_0^t b_s(X_s^{\mathbf y,\mathbf z},\mathbf y_s,\mathbf z_s)
      \,\mathrm ds\\
  &\quad+\int_0^t
    \sigma_s(X_s^{\mathbf y,\mathbf z},\mathbf y_s,\mathbf z_s)
    \,\mathrm dB_s,
\end{aligned}
\]
and then solve the Lipschitz BSDE
\[
\begin{aligned}
  Y_t^{\mathbf y,\mathbf z}
  &=g(X_T^{\mathbf y,\mathbf z})
    +\int_t^T f_s(
      X_s^{\mathbf y,\mathbf z},
      Y_s^{\mathbf y,\mathbf z},
      Z_s^{\mathbf y,\mathbf z}
    )\,\mathrm ds\\
  &\quad-\int_t^T Z_s^{\mathbf y,\mathbf z}\,\mathrm dB_s.
\end{aligned}
\]
Define
$F(\mathbf y,\mathbf z):=(Y^{\mathbf y,\mathbf z},
Z^{\mathbf y,\mathbf z})$ on the Banach space with norm
\[
  \|(\mathbf y,\mathbf z)\|_w^2
  :=\sup_{0\le t\le T}\mathbb E|\mathbf y_t|^2
    +\mathbb E\!\left[\int_0^T|\mathbf z_t|^2\,\mathrm dt\right].
\]

For two inputs, use $\Delta$ for differences.  Linearization gives
\[
\begin{aligned}
  \Delta X_t
  &=\int_0^t(\alpha_s^1\Delta X_s
    +\beta_s^1\Delta\mathbf y_s
    +\gamma_s^1\Delta\mathbf z_s)\,\mathrm ds\\
  &\quad+\int_0^t(\alpha_s^2\Delta X_s
    +\beta_s^2\Delta\mathbf y_s
    +\gamma_s^2\Delta\mathbf z_s)\,\mathrm dB_s,\\
  \Delta Y_t
  &=\lambda\Delta X_T
    +\int_t^T(\alpha_s^3\Delta X_s
      +\beta_s^3\Delta Y_s
      +\gamma_s^3\Delta Z_s)\,\mathrm ds
    -\int_t^T\Delta Z_s\,\mathrm dB_s,
\end{aligned}
\]
where the coefficients are bounded,
$|\gamma^2|\le L_{\sigma,z}$, and $|\lambda|\le L_g$.
For $0<\varepsilon<1/2$, It\^o's formula and Young's inequality yield,
provided $CT/\varepsilon\le\varepsilon$,
\[
  \sup_{t\le T}\mathbb E|\Delta X_t|^2
  \le\frac{L_{\sigma,z}^2+2\varepsilon}{1-\varepsilon}
      \|(\Delta\mathbf y,\Delta\mathbf z)\|_w^2
\]
and
\[
  \|(\Delta Y,\Delta Z)\|_w^2
  \le\frac{L_g^2+\varepsilon^2}{1-2\varepsilon}
      \sup_{t\le T}\mathbb E|\Delta X_t|^2.
\]
Consequently,
\[
  \|(\Delta Y,\Delta Z)\|_w^2
  \le c_\varepsilon
    \|(\Delta\mathbf y,\Delta\mathbf z)\|_w^2,
\]
where
\[
  c_\varepsilon
  :=\frac{(L_g^2+\varepsilon^2)(L_{\sigma,z}^2+2\varepsilon)}
          {(1-\varepsilon)(1-2\varepsilon)}.
\]
As $\varepsilon\downarrow0$,
$c_\varepsilon\to(L_gL_{\sigma,z})^2<1$.  Choose $\varepsilon$ so that
$c_\varepsilon<1$, and then choose $\delta_0$ so that
$CT/\varepsilon\le\varepsilon$ whenever $T\le\delta_0$.  Banach's fixed
point theorem gives a unique fixed point $(Y,Z)$; with
$X=X^{Y,Z}$, the triple $(X,Y,Z)$ is the unique solution of the FBSDE\@.

For the estimate, let $(Y^0,Z^0)=F(0,0)$.  The contraction property gives
\[
  \|(Y,Z)\|_w
  \le\frac{1}{1-\sqrt{c_\varepsilon}}
      \|(Y^0,Z^0)\|_w.
\]
Standard SDE and BSDE estimates give
$\|(Y^0,Z^0)\|_w\le C(|x|+I_0)$, and a final SDE estimate for $X$ proves
the asserted bound for $\Theta$.
\end{proof}

\begin{cor}[Local stability]
Let two data sets satisfy the hypotheses of the theorem with the same
local horizon.  Denote their solutions by $\Theta$ and
$\widetilde\Theta$, and write $\Delta\varphi=\widetilde\varphi-\varphi$.
Then
\[
\begin{aligned}
  \|\widetilde\Theta-\Theta\|^2
  \le C\mathbb E\!\Bigg[&|\widetilde x-x|^2
    +\left(\int_0^T
      \bigl(
        |\Delta b_t(\widetilde\Theta_t)|
        +|\Delta f_t(\widetilde\Theta_t)|
      \bigr)\,\mathrm dt\right)^2\\
    &+\int_0^T|\Delta\sigma_t(\widetilde\Theta_t)|^2\,\mathrm dt
    +|\Delta g(\widetilde X_T)|^2
  \Bigg].
\end{aligned}
\]
\end{cor}

\begin{remark}
The condition $L_{\sigma,z}L_g<1$ is sufficient, not necessary.  For
example,
\[
  X_t=x+\gamma\int_0^t Z_s\,\mathrm dB_s,
  \qquad
  Y_t=\lambda X_T-\int_t^T Z_s\,\mathrm dB_s
\]
is well posed if and only if $\gamma\lambda\ne1$.
\end{remark}

\section{The decoupling-field method}

On a local interval $[t,T]$, let $\Theta^{t,x}$ be the solution starting
from $X_t^{t,x}=x$, and define
\[
  u_t(x):=Y_t^{t,x}.
\]
For random coefficients, $u$ is generally a random field.

\begin{thm}[Local decoupling field]
Under the hypotheses of the small-time theorem:
\begin{itemize}
  \item[(i)] $u_t(x)\in L^2(\mathcal F_t;\mathbb R^{d_2})$ for every
  $(t,x)$;
  \item[(ii)] there is a deterministic constant $C$ such that, outside a
  single null set,
  \[
    |u_t(x_1)-u_t(x_2)|\le C|x_1-x_2|
  \]
  for all $t,x_1,x_2$;
  \item[(iii)] every local solution satisfies
  \[
    Y_s^{t,x}=u_s(X_s^{t,x}),
    \qquad t\le s\le T.
  \]
\end{itemize}
\end{thm}

\begin{proof}
Measurability follows from the construction of the local solution.  The
conditional version of the local stability estimate gives the pathwise
Lipschitz inequality.  Finally, uniqueness on the interval $[s,T]$
identifies the continuation of the solution from time $s$ with the
solution starting at $(s,X_s^{t,x})$, proving the flow identity.
\end{proof}

\begin{de}
An $\mathbb F$-measurable random field
\[
  u:[0,T]\times\mathbb R^{d_1}\times\Omega\longrightarrow\mathbb R^{d_2},
  \qquad u_T(x)=g(x),
\]
is a decoupling field if there exists $\delta>0$ such that, for every
$t_1<t_2$ with $t_2-t_1\le\delta$ and every
$\eta\in L^2(\mathcal F_{t_1};\mathbb R^{d_1})$, the FBSDE on
$[t_1,t_2]$ with $X_{t_1}=\eta$ and terminal condition
$Y_{t_2}=u_{t_2}(X_{t_2})$ has a unique solution satisfying
\[
  Y_t=u_t(X_t),
  \qquad t_1\le t\le t_2.
\]
The decoupling field is regular if it is uniformly Lipschitz continuous
in $x$.
\end{de}

\begin{remark}
A decoupling field is unique.  Indeed, on a sufficiently short interval,
start at an arbitrary deterministic state $x$ at its left endpoint.
Local uniqueness identifies the value of $u_t(x)$.  Repeating this
argument backward over a time partition identifies the field on the
whole interval.
\end{remark}

\begin{thm}[Global solution from a decoupling field]
If the FBSDE has a decoupling field $u$, then it has a unique global
solution, and
\[
  Y_t=u_t(X_t),
  \qquad 0\le t\le T.
\]
\end{thm}

\begin{proof}
Choose a partition
$0=t_0<\cdots<t_N=T$ with mesh at most the local length in the definition.
Starting with $X_0=x$, solve successively on $[t_{i-1},t_i]$ the FBSDE
with terminal condition $u_{t_i}(X_{t_i})$.  The identities
$Y_{t_i}=u_{t_i}(X_{t_i})$ make the pieces agree at the grid points, so
they paste into a global solution.  For uniqueness, restrict any global
solution to the final interval, use local uniqueness there, and proceed
backward interval by interval.
\end{proof}

\subsection{A structural condition for global regularity}

\begin{thm}\label{fbsdeDF}
Assume $d=d_1=d_2=1$, $\sigma$ is independent of $z$, the coefficients
are continuously differentiable in $(x,y,z)$ with bounded derivatives,
and there is a constant $c>0$ such that
\[
  \partial_y\sigma\,\partial_zb
  \le-c\left|
    \partial_yb
    +\partial_x\sigma\,\partial_zb
    +\partial_y\sigma\,\partial_zf
  \right|.
\]
Then the FBSDE has a regular decoupling field and hence a unique global
solution.
\end{thm}

\begin{proof}
Local well-posedness holds because $\partial_z\sigma=0$.  Consider a
local solution on $[t_0,t_1]$ with terminal condition
$Y_{t_1}=h(X_{t_1})$, where $h$ is Lipschitz.  Differentiate the FBSDE
with respect to the initial state and write
$\nabla\Theta=(\nabla X,\nabla Y,\nabla Z)$.  Before a possible zero of
$\nabla X$, define
\[
  \widehat Y:=\frac{\nabla Y}{\nabla X},
  \qquad
  \widehat Z:=\frac{\nabla Z}{\nabla X}
    -\widehat Y(\partial_x\sigma+\partial_y\sigma\,\widehat Y).
\]
It\^o's formula gives the scalar BSDE
\[
  \mathrm d\widehat Y_s
  =-\bigl(
    \partial_xf_s+a_s(\widehat Y_s)\widehat Y_s
    +\theta_s(\widehat Y_s)\widehat Z_s
  \bigr)\,\mathrm ds
  +\widehat Z_s\,\mathrm dB_s,
\]
with terminal value $h'(X_{t_1})$, where
\[
\begin{aligned}
  a(y)
  &:=\partial_yf+\partial_xb+\partial_zf\,\partial_x\sigma\\
  &\quad+
    (\partial_yb+\partial_zf\,\partial_y\sigma
      +\partial_zb\,\partial_x\sigma)y
    +\partial_zb\,\partial_y\sigma\,y^2,\\
  \theta(y)
  &:=\partial_zf+\partial_x\sigma
    +(\partial_zb+\partial_y\sigma)y.
\end{aligned}
\]
The forward variational equation can be written as a stochastic
exponential, which proves that $\nabla X$ never reaches zero.  The
structural assumption and completion of the square give a deterministic
upper bound
\[
  a(y)\le\overline L
  \qquad\text{for all }y,
\]
where $\overline L$ depends only on the derivative bounds and $c$.
First take a smooth terminal map $h$ and stop when
$|\widehat Y|+|\nabla X|+|\nabla X|^{-1}$ reaches a fixed level.  On the
stopped interval, $\theta(\widehat Y)$ is bounded, so Girsanov's theorem
and the linear representation of the characteristic BSDE give
\[
  |\widehat Y_t|
  \le \|h'\|_\infty e^{\overline L(t_1-t)}
      +L\int_t^{t_1}e^{\overline L(s-t)}\,\mathrm ds.
\]
The bound is independent of the stopping level; hence the stopping times
increase to $t_1$.  A smooth approximation of a merely Lipschitz terminal
map gives the same estimate with $\|h'\|_\infty$ replaced by the
Lipschitz constant of $h$.
Thus the Lipschitz constant of the decoupling field on a local interval
is controlled solely by the Lipschitz constant at the right endpoint and
fixed model constants.  Iterating this estimate backward from
$h=g$ gives a finite bound on the whole horizon.  The local existence
length can therefore be chosen uniformly, so the decoupling field
extends to time zero.  The preceding theorem then gives the global
solution.
\end{proof}

\section{Properties derived from the decoupling field}

\begin{thm}[Comparison of decoupling fields]
Let $(b,\sigma,f^i,g^i)$, $i=1,2$, satisfy the hypotheses of
Theorem~\ref{fbsdeDF}, and let $u^i$ be their regular decoupling fields.
If
\[
  f^1\le f^2
  \qquad\text{and}\qquad
  g^1\le g^2,
\]
then $u^1\le u^2$.
\end{thm}

\begin{proof}
It is enough to work on one common local interval and then iterate
backward.  Fix the same initial state for the two FBSDEs and linearize
the difference.  The terminal difference has the form
\[
  \Delta Y_T=\lambda\Delta X_T+\Delta g(X_T^1),
  \qquad \Delta g:=g^1-g^2\le0,
\]
and the generator difference is the sum of a linear term in
$(\Delta X,\Delta Y,\Delta Z)$ and
$\Delta f(\Theta^1)$, where $\Delta f=f^1-f^2\le0$.

Solve the associated homogeneous linear FBSDE with unit initial forward
state and terminal coefficient $\lambda$.  Its forward component is
strictly positive.  Dividing its backward component by its forward
component produces bounded processes $(\widehat Y,\widehat Z)$.  Set
\[
  \delta Y:=\Delta Y-\widehat Y\Delta X
\]
and choose the corresponding $\delta Z$ so that the coefficient of
$\Delta X$ cancels.  Then $(\delta Y,\delta Z)$ solves a scalar linear
BSDE with terminal value $\Delta g(X_T^1)\le0$ and inhomogeneous term
$\Delta f(\Theta^1)\le0$.  The linear representation formula, after the
standard BMO localization for its stochastic exponential, gives
$\delta Y_0\le0$.  Since $\Delta X_0=0$, this is
$u_0^1(x)-u_0^2(x)\le0$.  Applying the same argument from each starting
time and then iterating over the partition proves the result.
\end{proof}

\begin{remark}
This theorem compares the decoupling fields, not the two backward
processes along their different forward states.  In general,
$X_t^1\ne X_t^2$, so no pointwise comparison between
$Y_t^1=u_t^1(X_t^1)$ and $Y_t^2=u_t^2(X_t^2)$ follows for $t>0$.
\end{remark}

\begin{thm}[A variational estimate]
Consider the scalar linear FBSDE
\[
\left\{
\begin{aligned}
  X_t
  &=1+\int_0^t(\alpha_s^1X_s+\beta_s^1Y_s+\gamma_s^1Z_s)\,\mathrm ds
      +\int_0^t(\alpha_s^2X_s+\beta_s^2Y_s)\,\mathrm dB_s,\\
  Y_t
  &=GX_T+\int_t^T(\alpha_s^3X_s+\beta_s^3Y_s+\gamma_s^3Z_s)
      \,\mathrm ds
      -\int_t^T Z_s\,\mathrm dB_s.
\end{aligned}
\right.
\]
Assume
$|\alpha^i|,|\beta^i|,|\gamma^i|\le K$, $|G|\le K_0$, and
\[
  \beta^2\gamma^1=0,
  \qquad
  \beta^1+\alpha^2\gamma^1+\beta^2\gamma^3=0.
\]
For a sufficiently small horizon, the FBSDE has a unique solution and
\[
  |Y_t|\le\overline K_0|X_t|,
  \qquad
  \overline K_0:=(K_0+1)e^{(2K+K^2)T}-1.
\]
In particular, $|Y_0|\le\overline K_0$.
\end{thm}

\begin{proof}
Local well-posedness follows from the small-time theorem.  Conditional
stability, tested against arbitrary bounded $\mathcal F_t$-measurable
multipliers, first gives $|Y_t|\le C|X_t|$.  Before the first zero of
$X$, set
\[
  \widetilde Y:=Y/X,
  \qquad
  \widetilde Z:=Z/X-\widetilde Y(\alpha^2+\beta^2\widetilde Y).
\]
The two structural equalities eliminate the quadratic and linear powers
of $\widetilde Y$ from the drift coefficient generated by It\^o's
formula.  After localization at the levels $X=1/n$ and a Girsanov change
of measure, the linear representation gives
\[
  |\widetilde Y_t|
  \le K_0e^{(2K+K^2)(T-t)}
      +K\int_t^T e^{(2K+K^2)(s-t)}\,\mathrm ds
  \le\overline K_0.
\]
The stochastic-exponential representation of $X$ shows that it cannot
hit zero before the terminal time; the localization can therefore be
removed.  Multiplying by $|X_t|$ proves the estimate.
\end{proof}

\section{The method of continuation}

We now return to scalar FBSDEs and impose the monotonicity conditions:
there is $c>0$ such that, for every
$\theta_i=(x_i,y_i,z_i)$ and
$\Delta\theta=\theta_1-\theta_2$,
\[
\begin{aligned}
  &[b_t(\theta_1)-b_t(\theta_2)]\Delta y
   +[\sigma_t(\theta_1)-\sigma_t(\theta_2)]\Delta z\\
  &\qquad-[f_t(\theta_1)-f_t(\theta_2)]\Delta x
  \le-c(|\Delta x|^2+|\Delta y|^2+|\Delta z|^2),\\
  &[g(x_1)-g(x_2)]\Delta x\ge0.
\end{aligned}
\]

\begin{thm}[Uniqueness under monotonicity]
Under \textup{(F1)}--\textup{(F3)} and the monotonicity conditions, the
FBSDE has at most one solution.
\end{thm}

\begin{proof}
Let $\Theta^1$ and $\Theta^2$ be two solutions and use $\Delta$ for
differences.  It\^o's formula gives
\[
\begin{aligned}
  \mathrm d(\Delta X_t\Delta Y_t)
  &=\bigl(
    -\Delta f_t\Delta X_t
    +\Delta b_t\Delta Y_t
    +\Delta\sigma_t\Delta Z_t
  \bigr)\,\mathrm dt\\
  &\quad+\bigl(
    \Delta X_t\Delta Z_t
    +\Delta\sigma_t\Delta Y_t
  \bigr)\,\mathrm dB_t.
\end{aligned}
\]
Since $\Delta X_0=0$ and
$\Delta Y_T=g(X_T^1)-g(X_T^2)$, expectation and monotonicity yield
\[
  0\le
  -c\mathbb E\!\left[
    \int_0^T(|\Delta X_t|^2+|\Delta Y_t|^2+|\Delta Z_t|^2)\,\mathrm dt
  \right].
\]
Thus $\Delta Z=0$ almost everywhere and, by continuity of $X$ and $Y$,
$\Delta X=\Delta Y=0$ for every time.
\end{proof}

\begin{lem}[Solvable reference equation]
Let $b^0,f^0$ be progressively measurable with
\[
  \mathbb E\!\left[
    \left(\int_0^T(|b_t^0|+|f_t^0|)\,\mathrm dt\right)^2
  \right]<\infty,
\]
let $\sigma^0\in\mathbb H^2(0,T)$, and let
$g^0\in L^2(\mathcal F_T)$.  Then
\[
\left\{
\begin{aligned}
  X_t
  &=x+\int_0^t(-Y_s+b_s^0)\,\mathrm ds
      +\int_0^t(-Z_s+\sigma_s^0)\,\mathrm dB_s,\\
  Y_t
  &=X_T+g^0+\int_t^T(X_s+f_s^0)\,\mathrm ds
      -\int_t^T Z_s\,\mathrm dB_s
\end{aligned}
\right.
\]
has a unique solution.
\end{lem}

\begin{proof}
Uniqueness follows from the monotonicity estimate with $c=1$.  For
existence, set $\overline Y:=Y-X$.  Direct subtraction of the forward
and backward differentials shows that
\[
  \overline Y_t
  =g^0+\int_t^T(-\overline Y_s+f_s^0+b_s^0)\,\mathrm ds
      -\int_t^T\overline Z_s\,\mathrm dB_s,
\]
where
\[
  \overline Z=2Z-\sigma^0.
\]
This linear BSDE has a unique solution.  Set
\[
  Z:=\frac12(\overline Z+\sigma^0)
\]
and solve
\[
  X_t=x+\int_0^t(-X_s-\overline Y_s+b_s^0)\,\mathrm ds
      +\frac12\int_0^t(\sigma_s^0-\overline Z_s)\,\mathrm dB_s.
\]
Finally, define $Y=X+\overline Y$.  Substitution verifies the original
FBSDE\@.
\end{proof}

For $\alpha\in[0,1]$, define
\[
\begin{aligned}
  b_t^\alpha(x,y,z)&:=\alpha b_t(x,y,z)-(1-\alpha)y,\\
  \sigma_t^\alpha(x,y,z)&:=\alpha\sigma_t(x,y,z)-(1-\alpha)z,\\
  f_t^\alpha(x,y,z)&:=\alpha f_t(x,y,z)+(1-\alpha)x,\\
  g^\alpha(x)&:=\alpha g(x)+(1-\alpha)x.
\end{aligned}
\]
These coefficients satisfy the monotonicity condition with constant
\[
  c_\alpha:=\alpha c+1-\alpha\ge\min(c,1).
\]
Let $\operatorname{FBSDE}(\alpha)$ denote the equation with coefficients
$(b^\alpha,\sigma^\alpha,f^\alpha,g^\alpha)$ and arbitrary additive data
$(b^0,\sigma^0,f^0,g^0)$ in the spaces used in the preceding lemma.
Call $\operatorname{FBSDE}(\alpha)$ solvable if it has a solution for
every such additive data set.

\begin{lem}[Continuation step]
If $\operatorname{FBSDE}(\alpha_0)$ is solvable, then there is a number
$\delta_0>0$, depending only on the common Lipschitz and monotonicity
constants, such that
$\operatorname{FBSDE}(\alpha)$ is solvable for every
\[
  \alpha\in[\alpha_0,\min(1,\alpha_0+\delta_0)].
\]
\end{lem}

\begin{proof}
Write $\delta=\alpha-\alpha_0$.  Starting from $\Theta^0=0$, define
$\Theta^{n+1}$ as the solution of
$\operatorname{FBSDE}(\alpha_0)$ with additive data
\[
\begin{aligned}
  b_t^{n,0}&:=b_t^0+\delta\bigl(Y_t^n+b_t(\Theta_t^n)\bigr),\\
  \sigma_t^{n,0}&:=\sigma_t^0+\delta\bigl(Z_t^n+\sigma_t(\Theta_t^n)\bigr),\\
  f_t^{n,0}&:=f_t^0+\delta\bigl(f_t(\Theta_t^n)-X_t^n\bigr),\\
  g^{n,0}&:=g^0+\delta\bigl(g(X_T^n)-X_T^n\bigr).
\end{aligned}
\]
Solvability at $\alpha_0$ makes the iteration well defined.  Let
$\Delta\Theta^n=\Theta^{n+1}-\Theta^n$.  It\^o's formula for
$\Delta X^n\Delta Y^n$, the monotonicity of the
$\alpha_0$ coefficients, and Young's inequality control the natural
$L^2$ norm of $\Delta\Theta^n$ by the preceding increment.  The terminal
increment
\[
  \delta\bigl(
    g(X_T^n)-g(X_T^{n-1})-\Delta X_T^{n-1}
  \bigr)
\]
must be retained in this estimate; it is absorbed using the Lipschitz
bound for $g$ and the standard SDE estimate.  Altogether,
\[
  \|\Delta\Theta^n\|_{L^2(\mathrm dt\otimes\mathrm d\mathbb P)}^2
  \le C\delta^2
    \|\Delta\Theta^{n-1}\|_{L^2(\mathrm dt\otimes\mathrm d\mathbb P)}^2.
\]
The terms of order $\delta$ multiplying the current increment can be
absorbed into the monotonicity constant.  Choose $\delta_0$ so that
$C\delta_0^2<1/4$.  Then the iteration is Cauchy.  Passing to the limit
in the Lipschitz coefficients shows that its limit solves
$\operatorname{FBSDE}(\alpha_0+\delta)$.
\end{proof}

\begin{thm}[Global well-posedness by continuation]
Under \textup{(F1)}--\textup{(F3)} and the monotonicity conditions, the
FBSDE admits a unique solution on every finite horizon.
\end{thm}

\begin{proof}
The reference lemma proves solvability at $\alpha=0$.  Choose an integer
$N$ such that
\[
  (N-1)\delta_0<1\le N\delta_0.
\]
Applying the continuation-step lemma at most $N$ times reaches
$\alpha=1$, which is the original FBSDE\@.  Uniqueness was proved above.
\end{proof}
